\begin{pastexam}[label={px:ncku-114-ee-ds-03}]{成大電機所 114 學年度 單選題第 3 題}
（單選；每題 5 分，答錯倒扣 2 分）

Given a directed graph in the following, please identify which edge would become a cross edge if
we traverse the graph by the breadth first search algorithm from node C. Assume that the nodes
are stored in the adjacent list according to alphabetic order. Please specify which edge forms a
cross edge in the resulting breadth-first spanning tree.

\begin{center}
\begin{tikzpicture}[v/.style={draw,circle,minimum size=7mm,inner sep=0pt},>={Stealth[length=2.2mm]}]
  \node[v] (A) at (0,0) {A};
  \node[v] (B) at (1.7,0.9) {B};
  \node[v] (C) at (2.1,-0.8) {C};
  \node[v] (D) at (4.1,0) {D};
  \node[v] (E) at (6.0,0) {E};
  \node[v] (F) at (6.0,1.5) {F};
  \draw[->] (A) -- (B);
  \draw[->] ([yshift=1.5mm]C.west) -- ([yshift=-0.5mm]A.east);
  \draw[->] (A.south east) -- ([yshift=-1.5mm]C.west);
  \draw[->] ([xshift=-1mm]B.south) -- ([xshift=-2.5mm]C.north);
  \draw[->] ([xshift=1mm]C.north) -- ([xshift=2.5mm]B.south);
  \draw[->] (B.east) -- (D.north west);
  \draw[->] ([yshift=2.5mm]D.north west) -- ([yshift=2.5mm]B.east);
  \draw[->] (C) -- (D);
  \draw[->] (D) -- (F);
  \draw[->] (D) -- (E);
  \draw[->] (F) -- (E);
\end{tikzpicture}
\end{center}
\begin{enumerate}[label=\alph*)]
\item (F, E)
\item (B, D)
\item (D, B)
\item (A, B)
\item All of the above.
\end{enumerate}
\end{pastexam}

\begin{pastsoln}
答案是 e) All of the above。

圖的邊：A$\to$B、A$\to$C、B$\to$C、B$\to$D、C$\to$A、C$\to$B、C$\to$D、D$\to$B、D$\to$E、
D$\to$F、F$\to$E。 鄰接串列照字母排：A: B, C；B: C, D；C: A, B, D；D: B, E, F；E: （空）；F: E。

\textbf{BFS 從 C 開始}（資結書第~21.3~節）：
\begin{center}
\begin{tabular}{clll}
\toprule
取出 & 檢查的邊 & 新發現（樹邊） & 佇列 \\
\midrule
C & C$\to$A, C$\to$B, C$\to$D & A, B, D（層 1） & A, B, D \\
A & A$\to$B（B 已發現）, A$\to$C（C 已發現） & & B, D \\
B & B$\to$C, B$\to$D（都已發現） & & D \\
D & D$\to$B（已發現）, D$\to$E, D$\to$F & E, F（層 2） & E, F \\
E & （沒有出邊） & & F \\
F & F$\to$E（已發現） & & （空） \\
\bottomrule
\end{tabular}
\end{center}
BFS 樹：C 是根，子節點 A、B、D；D 的子節點 E、F。

\textbf{非樹邊分類}（資結書第~21.5~節）：指向祖先的是後向邊，兩端沒有祖孫關係的是交叉邊。
\begin{itemize}
\item A$\to$C、B$\to$C：C 是 A、B 的祖先（根），後向邊。
\item A$\to$B：A、B 都是 C 的子節點，互不為祖孫，\textbf{交叉邊}。
\item B$\to$D、D$\to$B：B、D 是兄弟，\textbf{交叉邊}。
\item F$\to$E：E、F 都是 D 的子節點，\textbf{交叉邊}。
\end{itemize}
有向圖的 BFS 不會有前向邊（若 $u$ 有邊到後代 $v$，$v$ 早就在處理 $u$ 時被發現成為 $u$ 的子節點）。
選項 a)--d) 的四條邊都是交叉邊，所以選 e)。
\end{pastsoln}
